% GENERATED FILE. Edit canonical lessons, then run npm run generate:books.
% canonical-source-hash: fnv1a64:b65230f10ead698b
% canonical-lessons: GE-C416-einchecken, GE-C416-bordkarte, GE-C416-flugsteig, GE-C416-handgepaeck, GE-C416-aufgeben

\chapter{Am Flughafen}
\label{ch:ge-am-flughafen}
\hlchaptermodality{416}

\begin{chapteropening}
\textbf{By the end of this chapter:} \emph{I can check in for a flight: checking in, the boarding pass, the gate, hand luggage and checking a bag in.}

Complete the last lesson of chapter 416: I can check in for a flight: checking in, the boarding pass, the gate, hand luggage and checking a bag in.
\end{chapteropening}

\section[einchecken]{einchecken — to check in}

\label{lesson:GE-C416-einchecken}



\begin{quote}
Before the new one: say the German for to ride without a ticket, then the German for to queue.
\end{quote}

\subsection*{You'll want to know: einchecken}
\textbf{einchecken} — \textquotedblleft{}to check in\textquotedblright{}.

\begin{quote}
Zwei Stunden vor dem Abflug checken wir ein.
\end{quote}

\begin{quote}
\emph{Two hours before departure we check in.}
\end{quote}

Separable: \emph{ich checke ein}, \emph{ich habe eingecheckt}. Also at a hotel.

\begin{cousinweb}
\emph{einchecken} is built from \emph{ein-} 'in' and \emph{checken}, which was borrowed from English \emph{check} (EtymDB), on the model of English \emph{check in}.
\end{cousinweb}

\begin{grammarlens}[title={What you've built}]
One more word for telling, explaining and describing at B1.
\end{grammarlens}

\subsection*{Your turn}
\begin{itemize}
\raggedright
  \item \emph{Say it:} \emph{einchecken}
  \item \emph{Say it:} \emph{einchecken}, once more
  \item \emph{Say it:} a sentence of your own with \emph{einchecken}
  \item \emph{Recall:} say the German for to ride without a ticket, then the German for to queue, then say \emph{einchecken} again
\end{itemize}

\begin{tcolorbox}[breakable,colback=teal!4,colframe=teal!35!black,title={Before you move on}]
What does einchecken mean? (\textquotedblleft{}To check in\textquotedblright{}.) Say it once more.
\end{tcolorbox}
\section[die Bordkarte]{die Bordkarte — the boarding pass}

\label{lesson:GE-C416-bordkarte}



\begin{quote}
Before the new one: say the German for to queue, then the German for to check in.
\end{quote}

\subsection*{You'll want to know: die Bordkarte}
\textbf{die Bordkarte} — \textquotedblleft{}the boarding pass\textquotedblright{}.

\begin{quote}
Die Bordkarte ist auf ihrem Handy.
\end{quote}

\begin{quote}
\emph{The boarding pass is on her phone.}
\end{quote}

Plural: \emph{die Bordkarten}. \emph{An Bord} — on board.

\begin{cousinweb}
\emph{Bordkarte} is built from \emph{Bord} 'board, side of a ship', inherited from Old High German \emph{bort} 'edge, rim' (EtymDB), and \emph{Karte} 'card', which goes back to Latin \emph{charta} 'sheet of paper' (EtymDB; Lewis \& Short).
\end{cousinweb}

\begin{grammarlens}[title={What you've built}]
One more word for telling, explaining and describing at B1.
\end{grammarlens}

\subsection*{Your turn}
\begin{itemize}
\raggedright
  \item \emph{Say it:} \emph{die Bordkarte}
  \item \emph{Say it:} \emph{die Bordkarte}, once more
  \item \emph{Say it:} a sentence of your own with \emph{die Bordkarte}
  \item \emph{Recall:} say the German for to queue, then the German for to check in, then say \emph{die Bordkarte} again
\end{itemize}

\begin{tcolorbox}[breakable,colback=teal!4,colframe=teal!35!black,title={Before you move on}]
What does die Bordkarte mean? (\textquotedblleft{}The boarding pass\textquotedblright{}.) Say it once more.
\end{tcolorbox}
\section[der Flugsteig]{der Flugsteig — the gate (at an airport)}

\label{lesson:GE-C416-flugsteig}



\begin{quote}
Before the new one: say the German for to check in, then the German for the boarding pass.
\end{quote}

\subsection*{You'll want to know: der Flugsteig}
\textbf{der Flugsteig} — \textquotedblleft{}the gate (at an airport)\textquotedblright{}.

\begin{quote}
Der Flugsteig für unseren Flug ist ganz am Ende.
\end{quote}

\begin{quote}
\emph{The gate for our flight is right at the end.}
\end{quote}

Plural: \emph{die Flugsteige}. People also say \emph{das Gate}.

\begin{cousinweb}
\emph{Flugsteig} is built from \emph{Flug} 'flight', inherited from Proto-Germanic \emph{flugiz} (EtymDB), and \emph{Steig} 'path, way up', from \emph{steigen} 'to climb'.
\end{cousinweb}

\begin{grammarlens}[title={What you've built}]
One more word for telling, explaining and describing at B1.
\end{grammarlens}

\subsection*{Your turn}
\begin{itemize}
\raggedright
  \item \emph{Say it:} \emph{der Flugsteig}
  \item \emph{Say it:} \emph{der Flugsteig}, once more
  \item \emph{Say it:} a sentence of your own with \emph{der Flugsteig}
  \item \emph{Recall:} say the German for to check in, then the German for the boarding pass, then say \emph{der Flugsteig} again
\end{itemize}

\begin{tcolorbox}[breakable,colback=teal!4,colframe=teal!35!black,title={Before you move on}]
What does der Flugsteig mean? (\textquotedblleft{}The gate (at an airport)\textquotedblright{}.) Say it once more.
\end{tcolorbox}
\section[das Handgepäck]{das Handgepäck — the hand luggage, the carry-on}

\label{lesson:GE-C416-handgepaeck}



\begin{quote}
Before the new one: say the German for the boarding pass, then the German for the gate.
\end{quote}

\subsection*{You'll want to know: das Handgepäck}
\textbf{das Handgepäck} — \textquotedblleft{}the hand luggage, the carry-on\textquotedblright{}.

\begin{quote}
Im Handgepäck hat er nur ein Buch und eine Jacke.
\end{quote}

\begin{quote}
\emph{In his hand luggage he has only a book and a jacket.}
\end{quote}

No plural.

\begin{cousinweb}
\emph{Handgepäck} is built from \emph{Hand}, inherited from Proto-Germanic \emph{handuz}, and \emph{Gepäck} 'luggage', from \emph{packen} 'to pack'.
\end{cousinweb}

\begin{grammarlens}[title={What you've built}]
One more word for telling, explaining and describing at B1.
\end{grammarlens}

\subsection*{Your turn}
\begin{itemize}
\raggedright
  \item \emph{Say it:} \emph{das Handgepäck}
  \item \emph{Say it:} \emph{das Handgepäck}, once more
  \item \emph{Say it:} a sentence of your own with \emph{das Handgepäck}
  \item \emph{Recall:} say the German for the boarding pass, then the German for the gate, then say \emph{das Handgepäck} again
\end{itemize}

\begin{tcolorbox}[breakable,colback=teal!4,colframe=teal!35!black,title={Before you move on}]
What does das Handgepäck mean? (\textquotedblleft{}The hand luggage, the carry-on\textquotedblright{}.) Say it once more.
\end{tcolorbox}
\section[aufgeben]{aufgeben — to give up; to check in (a bag); to post}

\label{lesson:GE-C416-aufgeben}



\begin{quote}
Before the new one: say the German for the gate, then the German for the hand luggage.
\end{quote}

\subsection*{You'll want to know: aufgeben}
\textbf{aufgeben} — \textquotedblleft{}to give up; to check in (a bag); to post\textquotedblright{}.

\begin{quote}
Den großen Koffer gibt sie auf, den kleinen nimmt sie mit.
\end{quote}

\begin{quote}
\emph{She checks in the big suitcase and takes the small one with her.}
\end{quote}

Separable, strong: \emph{gibt auf}, \emph{gab auf}, \emph{hat aufgegeben}. The everyday sense is to give up: \emph{Er gibt nie auf.}

\begin{cousinweb}
\emph{aufgeben} is built from \emph{auf} 'up' and \emph{geben} 'to give', inherited from Proto-Germanic \emph{gebaną} (EtymDB). EtymDB pairs it with the rare English \emph{upgive}.
\end{cousinweb}

\begin{grammarlens}[title={What you've built}]
One more word for telling, explaining and describing at B1.
\end{grammarlens}

\subsection*{Your turn}
\begin{itemize}
\raggedright
  \item \emph{Say it:} \emph{aufgeben}
  \item \emph{Say it:} \emph{aufgeben}, once more
  \item \emph{Say it:} a sentence of your own with \emph{aufgeben}
  \item \emph{Recall:} say the German for the gate, then the German for the hand luggage, then say \emph{aufgeben} again
\end{itemize}

\begin{tcolorbox}[breakable,colback=teal!4,colframe=teal!35!black,title={Before you move on}]
What does aufgeben mean? (\textquotedblleft{}To give up; to check in (a bag); to post\textquotedblright{}.) Say it once more.
\end{tcolorbox}
